\section{Graph Model} \label{sec:graph-model}
In this section, we introduce a graph representation of the core genomes that
enables us to efficiently solve the inference problem introduced in \cref{sec:alignment}.
The third column in \cref{fig:models} shows the resulting generative model.

Variation graphs provide a succinct encoding of all possible segments of core genomes.
Conventional variation graphs have shown great promise in capturing large-scale variations that linear references cannot represent \cite{VGtool,paten2017genome,garrison2018variation}.
Here, we generalize these variation graphs to our probabilistic setting.

\begin{figure*}
	\begin{minipage}{0.7\linewidth}
		\begin{tikzpicture}[
		node distance=0.5cm,
		myletter/.style={letter,scale=0.9,text width=2em,inner sep=0pt,align=center},
		myarrow/.style={arrow}
	]
	% vertex positions
	\def \vd {0.8}
	\pgfmathsetmacro{\leftx}{3}
	\pgfmathsetmacro{\middlex}{\leftx+4*\vd}
	\pgfmathsetmacro{\rightx}{\leftx+7*\vd}

	% GENOMES

	\node[] at (0, 0.9) (g1) {$\genome_1^\gcore=$ \textcolor{my-full-blue}{$AAC$} $CCG$ \textcolor{my-full-blue}{$AAC$}};
	\node[] at (0, 0.3) (g2) {$\genome_2^\gcore=$ \textcolor{my-full-blue}{$AAC$} $CCG$ \textcolor{my-full-purple}{$GGT$}};
	\node[] at (0,-0.3) (g3) {$\genome_3^\gcore=$ \textcolor{my-full-purple}{$GGT$} $CCG$ \textcolor{my-full-blue}{$AAC$}};
	\node[] at (0,-0.9) (g4) {$\genome_4^\gcore=$ \textcolor{my-full-purple}{$GGT$} $CCG$ \textcolor{my-full-purple}{$GGT$}};

	% ARROW

	\node[inner sep=0pt,above right,anchor=center] at (2.2,0) {\includegraphics[scale=1]{figures/arrow}};

	% GRAPH-LEFT

	\node[myletter] at (\leftx, 0) (left) {$v_0$};

	\path (left)
		 ++(\vd,0.6) node[myletter,my-full-blue] (left-AAC-1) {$v_1$}
		 ++(\vd,0) node[myletter,my-full-blue] (left-AAC-2) {$v_2$}
		 ++(\vd,0) node[myletter,my-full-blue] (left-AAC-3) {$v_3$};
	
	\path (left)
		 ++(\vd,-0.6) node[myletter,my-full-purple] (left-GGT-1) {$v_4$}
		 ++(\vd,0) node[myletter,my-full-purple] (left-GGT-2) {$v_5$}
		 ++(\vd,0) node[myletter,my-full-purple] (left-GGT-3) {$v_6$};

	% GRAPH-MIDDLE

	\path (\middlex,0) node[myletter] (CCG-1) {$v_7$}
		++(\vd,0) node[myletter] (CCG-2) {$v_8$}
		++(\vd,0) node[myletter] (CCG-3) {$v_9$};
	
	% GRAPH-RIGHT

	\path (\rightx, 0.6) node[myletter,my-full-blue] (right-AAC-1) {$v_{10}$}
		   ++(\vd,0) node[myletter,my-full-blue] (right-AAC-2) {$v_{11}$}
		   ++(\vd,0) node[myletter,my-full-blue] (right-AAC-3) {$v_{12}$};

	\path (\rightx,-0.6) node[myletter,my-full-purple] (right-GGT-1) {$v_{13}$}
		++(\vd,0) node[myletter,my-full-purple] (right-GGT-2) {$v_{14}$}
		++(\vd,0) node[myletter,my-full-purple] (right-GGT-3) {$v_{15}$};
	
	% GRAPH EDGES

	\draw[myarrow,my-full-blue] (left) -- node[above] {$A$} (left-AAC-1);
	\draw[myarrow,my-full-blue] (left-AAC-1) -- node[above] {$A$} (left-AAC-2);
	\draw[myarrow,my-full-blue] (left-AAC-2) -- node[above] {$C$} (left-AAC-3);

	\draw[myarrow,my-full-purple] (left) -- node[above] {$G$} (left-GGT-1);
	\draw[myarrow,my-full-purple] (left-GGT-1) -- node[above] {$G$} (left-GGT-2);
	\draw[myarrow,my-full-purple] (left-GGT-2) -- node[above] {$T$} (left-GGT-3);

	\draw[myarrow] (left-AAC-3) -- node[right,pos=0.3] {$C$} (CCG-1);
	\draw[myarrow] (left-GGT-3) -- node[right,pos=0.3] {$C$} (CCG-1);
	\draw[myarrow] (CCG-1) -- node[above] {$C$} (CCG-2);
	\draw[myarrow] (CCG-2) -- node[above] {$G$} (CCG-3);

	\draw[myarrow,my-full-blue] (CCG-3) -- node[left,pos=0.7] {$A$} (right-AAC-1);
	\draw[myarrow,my-full-blue] (right-AAC-1) -- node[above] {$A$} (right-AAC-2);
	\draw[myarrow,my-full-blue] (right-AAC-2) -- node[above] {$C$} (right-AAC-3);

	\draw[myarrow,my-full-purple] (CCG-3) -- node[left,pos=0.7] {$G$} (right-GGT-1);
	\draw[myarrow,my-full-purple] (right-GGT-1) -- node[above] {$G$} (right-GGT-2);
	\draw[myarrow,my-full-purple] (right-GGT-2) -- node[above] {$T$} (right-GGT-3);
\end{tikzpicture}
		\caption{Variation graph capturing genomes $\{\genome_1, \genome_2, \genome_3, \genome_4\}$.}
		\label{fig:deterministic-variation-graph}
	\end{minipage}\hspace{0.04\linewidth}
	\begin{minipage}{0.25\linewidth}
		\centering
		\includegraphics[width=0.8\linewidth]{figures/partition}
		\caption{Concretization of paths of length $L$.}
		\label{fig:partition}
	\end{minipage}
\end{figure*}

\para{Variation Graphs}
\cref{fig:deterministic-variation-graph} shows a set of genomes $\genomes$ (left) encoded by a conventional variation graph $\graph$ (right).
The graph $\graph=(V,E)$ consists of vertices $V$ and edges $E$ and is equipped with edge labels $\labels \colon E \to \{A,C,G,T\}$.
Thus, each path $\graphpath=e_1,\dots, e_n$ in $\graph$ spells a genome segment by $\labels(\graphpath)=\labels(e_1) \cdots \labels(e_n)$, where $\cdot$ denotes string concatenation.

In addition, $\graph$ is equipped with a position map $\position \colon E \to \mathcal{P}(\genomes \times \mathbb{N})$ which annotates each edge with a set of positions in the genomes.
For example, $\position((v_0,v_1))=\{\genome_1\llbracket 0 \rrbracket,\genome_2\llbracket 0 \rrbracket\}$, as the edge between $v_0$ and $v_1$ represents the first letter of the first two genomes.
In practice, we may not know the positions referred to by each vertex $v$.
In this case, we simply annotate $v$ with basic information about its position,
such as its approximate location in a reference genome.

When given a read $\segment^\gread$, existing graph alignment tools aim to find a path $\graphpath$ in $\graph$ that spells $\segment^\gread$, or a slightly edited version of it.
The annotations on the graph's vertices then allow us to identify the (set of) positions $\segment^\gread$ most likely originates from.

\para{Concrezitation of Paths}
Every path $\graphpath$ in $\graph$ represents one or multiple segments of genomes.
We capture these segments by a concretization function $\gamma$ taking as argument a path $\graphpath$ consisting of edges $e_i$.
Given $\graphpath$, $\gamma$ returns the set of all possible segments $\graphpath$ may represent, where each segment is encoded as a sequence of positions, determined by the position map $m$:
\begin{align*}
	\gamma(e_1, \dots, e_L) = \{p_1, \dots, p_L \mid p_i \in \position(e_i)\}
\end{align*}

\para{Compression by Variantion Graphs}
\cref{fig:partition} visualizes the concretization of paths of length $L$.
Green circles (\coloredcircle{my-full-green}) represent segments present in $\genomes$, while orange circles (\coloredcircle{my-full-orange}) represent spurious segments.
\cref{fig:partition} shows that the concretization of every path of length $L$ contains at least one segment from $\genomes$.
This is necessary when aligning reads of length $L$ since otherwise, alignment on $\graph$ may select a spurious path that does not reflect an existing genome segment.
Likewise, \cref{fig:partition} does not show any genome segment of length $L$ that is not captured by concretizing some path.
Otherwise, some reads of length $L$ could never be aligned correctly in $\graph$.

When $\graph$ captures these two properties for every length $L'<L$, we say $\graph$ \emph{non-deterministically $L$-compresses} the set of genomes $\genomes$.
This ensures that (i)~any path selected by alignment can be interpreted in the genomes and (ii)~any read can be aligned on $\graph$.
The formal definition is that for all lengths $L' \leq L$:
\begin{align}
	\forall \graphpath \in \paths{L'}{\graph}. \exists \genome \in \genomes, \startInGenome < \abs{\genome}-L'. \genome \llbracket \startInGenome \shortcolon \startInGenome + L' \rrbracket \in \gamma(\graphpath) \label{eq:compress-non-deterministic-precise} \\
	\forall \genome \in \genomes, \startInGenome < \abs{\genome}-L'. \exists \graphpath \in \paths{L'}{\graph}. \genome \llbracket \startInGenome \shortcolon \startInGenome + L' \rrbracket \in \gamma(\graphpath) \label{eq:compress-non-deterministic-complete}
\end{align}
% \begin{align} \label{eq:compress-non-deterministic}
% \bigcup_{\pi \in \graph} \gamma(\pi) = \left\{ \genome^\gcore\llbracket i \shortcolon i+L' \rrbracket \,\middle|\,
% \genome^\gcore \in \genomes^\gcore,
% 0 \leq i < |\genome^\gcore|-L'
% \right\}.
% \end{align}

% Generally, when aligning reads of length $L$ to a graph $\graph$, we need $\graph$ to $L$-compress the underlying set of core genomes.
% If \cref{eq:compress-non-deterministic-complete} does not hold, \ie $\graph$ does not represent some core genome segments, then some reads cannot be correctly aligned.
% If \cref{eq:compress-non-deterministic-precise} does not hold, \ie $\graph$ contains genome segments that do not occur in the core genomes, alignment on $\graph$ may select a spurious path that was only introduced in the graph model and does not reflect an existing core genome.

\para{Probabilistic Variation Graphs}
In order to account for our generative model, we generalize variation graphs to a probabilistic setting.
To this end, we extend $\graph$ by transition probabilities $\trans \colon E \to [0,1]$ which must induce a probability distribution over all outgoing edges, \ie for all $v \in V$, $\sum_{(v,v') \in E} \trans(v,v')=1$.
Then, the probability of selecting $\graphpath$ is given by $\trans(\graphpath)=\prod_{i=1}^{n} \trans(e_i)$.

In addition, we mark a node $\startNode \in V$ as a distinguished starting node.
Then, all paths $\graphpath=e_1,\dots,e_n$ in $\graph$ must start from $\startNode$, \ie $e_1=(\startNode,v)$ for some $v$.

\para{Compression by Probabilistic Variation Graphs}
To enable alignment on probabilistic variation graphs, it does not suffice if $\graph$ non-deterministically $L$ compresses $\genomes$.
Instead, we must additionally ensure that $\graph$ correctly captures the distribution of $\genomeRV$.
For example, in \cref{fig:partition}, we require that
$$\trans(\pi_3) = \Pr{\genomeRV \llbracket \startInGenomeRV \shortcolon \startInGenomeRV+L' \rrbracket \in \{\coloredcircle[1]{my-full-green},\coloredcircle[2]{my-full-green},\coloredcircle[3]{my-full-orange}\}}.$$

Formally, a probabilistic variation graph $\graph$ \emph{(probabilistically) $L$-compresses} a distribution over genomes $\genomeRV$ if the probability of every path $\graphpath$ of length $L' \leq L$ in $\graph$ matches the probability of the concretization of this path:
\begin{align} \label{eq:compress-probabilistic}
	\trans(\graphpath) = \Pr{\genomeRV\llbracket \startInGenomeRV \shortcolon \startInGenomeRV+L' \rrbracket \in \gamma(\graphpath)} && \forall \graphpath \in \paths{L'}{\graph}
\end{align}

\cref{lem:l-compress-generalize} states that probabilistic L-compression indeed
generalizes deterministic L-compression to the probabilistic setting.

\begin{lemma} \label{lem:l-compress-generalize}
	Assume $\graph$ does not contain edges that may never be used, \ie assume $\forall e \in E. \trans(e)>0$. If $\graph$ probabilistically $L$-compresses a distribution over genomes $\genomeRV$, then $\graph$ non-deterministically $L$-compresses the underlying set of genomes $\genomes$.
\end{lemma}
\begin{proof}
	Let $L' < L$.

	Showing \cref{eq:compress-non-deterministic-precise}: Let $\graphpath = e_1,\dots,e_n \in \paths{L'}{\graph}$. Because $\trans(\graphpath) > 0$ (otherwise there would be an edge that may never be used), and due to \cref{eq:compress-probabilistic}, $\Pr{\genomeRV \llbracket \startInGenomeRV \shortcolon \startInGenomeRV+L' \rrbracket \in \gamma(\graphpath)}>0$. Therefore, there must be at least one $\genome$, $\startInGenome$ with $\genomeRV\llbracket \startInGenomeRV \shortcolon \startInGenomeRV+L' \rrbracket \in \gamma(\graphpath)$.

	Showing \cref{eq:compress-non-deterministic-complete}: Due to \cref{eq:compress-probabilistic}, and because the transition probabilities $\trans$ induce a probability distribution, we have
	\begin{align*}
		\sum_{\pi \in \paths{L'}{\graph}} \Pr{\genomeRV\llbracket \startInGenomeRV \shortcolon \startInGenomeRV+L' \rrbracket \in \gamma(\graphpath)} = \sum_{\pi \in \paths{L'}{\graph}} \trans(\pi) = 1.
	\end{align*}
	Therefore, we know that all genome sequences must be covered by paths in $\graph$.
\end{proof}

\para{Approximate Compression}
In practice, we can only approximately satisfy \cref{eq:compress-probabilistic}.
\todo{Extend probabilistic L-compression to its approximate version}

% We measure the error according to KL divergence. 
% Concretely, we say that $\graph$ \emph{$(L,\epsilon)$-compresses} a distribution over genomes $\genomeRV$ if for all $L' \leq L$: \todo{explain more, if we are to stick with KL divergence}
% \begin{align*}
% 	\E[\Pi \sim \paths{L'}{\graph}]{
% 		\frac{
% 			\trans(\Pi)
% 		}{
% 			\Pr{\genomeRV\llbracket \startInGenomeRV \shortcolon \startInGenomeRV+L'\rrbracket \in \gamma(\Pi)}
% 		}
% 	} \leq \epsilon
% \end{align*}
% This essentially limits KL divergence: $\kld{\graph}{\genomeRV}\leq \epsilon$ (modulo $L$).

\subsection{Generative Graph Model}
\begin{figure}
	\lstinputlisting[language=MyPython,caption={Code for generating a read using the graph model.},label={lst:graph-minified}]{code/graph_minified.py}
\end{figure}

In \cref{lst:graph-minified}, we show an algorithmic description of our
generative model. It starts from an empty read (\cref{line:empty}) and extends
it until it reaches $L$ letters (\cref{line:empty}). In each iteration, it
samples a new edge in the graph (\cref{line:edge}), possibly mutates the letter
associated with this edge (\cref{line:mutate}) and records the resulting action
(\cref{line:record}). Unless the mutation (\cref{line:mutate}) results in a
deletion, the model reads the resulting letter (\cref{line:read}) possibly
resulting in a read error, and appends the resulting letter to the read. In
addition, unless the mutation (\cref{line:mutate}) results in an insertion (in
which case the next letter to be read remains the same), \cref{line:new-state}
updates the current position in the graph.

\subsection{Equivalence of Graph Model and Core Genomes Model}
\todo{Discuss and prove}


\section{Old Graph Model}

Instead of selecting an entire core genome $\genome^\gcore$, we generate a core segment $\segment^\gcore$ of $L$ letters that follows the same distribution as (i)~sampling a core genome $\genome^\gcore$ and (ii)~emitting $L$ letters from a random starting point in $\genome^\gcore$.

To select the core sequence $\seq^\gcore$, we sample from a variation graph, which provides a succinct representation of genomes in $\genomes^\gcore$.
A variation graph $H=(V,E)$ consists of vertices $V$ and directed edges $E \subseteq V \times V$ and is equipped with edge labels $l \colon E \to \{A,C,G,T\}$, transition probabilities $p \colon E \to [0,1]$, and a distinguished starting node $\startNode \in V$.
%
We assume that the probabilities $p$ induce a probability distribution over all outgoing edges, \ie for all $v \in V$, $\sum_{(v,v') \in E} p(v,v')=1$.
Starting from $\startNode$, we can sample a random walk $p=e_1,\dots,e_n$ in $H$ (where $e_1=(\startNode,v)$ for some $v$), whose probability is given by $\prod_{i=1}^{n} p(e_i)$.
Every path $p$ spells part of a reference genome by $l(e_1) \cdots l(e_n)$, where $\cdot$ denotes string concatenation.
%
\cref{fig:construct-variance-graph} shows a general construction, which builds a variance graph $H$ from a set of core genomes $\genomes^\gcore$ by $|\genomes^\gcore|$ separate paths, each spelling one core genome, and a starting node that can jump to any location of any genome in $\genomes^\gcore$, as long as this location still allows spelling a sequence of length $L$.
In practice, we will aim for a more compact representation of $H$. \todo{reference graph building section}

\begin{figure}
	\begin{tikzpicture}
		\def \len {6}
		\def \n {1}
		\def \jump {4}
		\node[letter] at (-0.5,0) (s) {$\startNode$};
		\foreach \gen in {-\n,...,\n} {
			\pgfmathtruncatemacro{\num}{\n-\gen}
			% nodes
			\foreach \pos in {1,...,\len} {
				\node[letter] at (\pos*0.7,\gen) (genome-\pos-\gen) {};
			}
			\foreach \pos in {1,...,\len} {
				% path edges
				\ifnum \numexpr\pos > 1
					\pgfmathtruncatemacro{\prev}{\pos-1}
					\ifnum \numexpr\gen = -1
						\draw[->] (genome-\prev-\gen) edge node[above] {$1$} (genome-\pos-\gen);
					\else
						\draw[->] (genome-\prev-\gen) edge node[below] {$1$} (genome-\pos-\gen);
					\fi
				\fi
				% jump edges
				\ifnum \numexpr\pos < \numexpr\jump+1
					\ifthenelse{\pos = 1}{
						\pgfmathsetmacro{\in}{180}
					}{
						\pgfmathsetmacro{\in}{180-\gen*30}
					}
					\ifthenelse{\pos = \jump}{
						\def \prob {$\frac{p_{\num}}{\jump}$}
					}{
						\def \prob {}
					}
					\ifnum \numexpr\gen = 1
						\draw[->] (s) edge[out=60+\pos*7,in=\in] node[above left,pos=0.2] {\prob} (genome-\pos-\gen);
					\fi
					\ifnum \numexpr\gen = 0
						\draw[->] (s) edge[out=10+\pos*5,in=180-30] node[above,pos=0.2] {\prob} (genome-\pos-\gen);
					\fi
					\ifnum \numexpr\gen = -1
						\draw[->] (s) edge[out=-60-\pos*7,in=\in] node[below left,pos=0.2] {\prob} (genome-\pos-\gen);
					\fi
				\fi
			}
			\node[right of=genome-\len-\gen] {$g_{\num}$};
		}
	\end{tikzpicture}
	\caption{General construction of a variance graph from core genomes. Each outgoing edge of $\startNode$ jumping into the \nth{i} genome $g_i$ has transition probability $\frac{p_i}{4}$. We do not allow jumping to the end of the genomes, as we must be able to read sequences of length $L$ (here, we depict $L=3$).}
	\label{fig:construct-variance-graph}
\end{figure}


To mutate the core sequence $\seq^\gcore$, we assume a process that copies $\seq^\gcore$ letter by letter, generating a mutated sequence $\seq^\gmut$.
When the process is about to copy the \nth{i} letter of $\seq^\gcore$, it probabilistically decides to (i)~insert a new letter (insertion; with probability $\pins$), (ii)~skip this letter (deletion; $\pdel$), (iii)~copy the letter incorrectly (edit; $\ped$), or (iv)~copy the letter correctly (copy; $\pcopy=1-\pins-\pdel-\ped$).
We assume that the new letter for edits is selected uniformly at random from all letters different from the current one, \eg copying letter $A$ yields (edited) letter $C$, $G$, or $T$ with a probability of $\frac{\ped}{3}$, respectively.

Finally, to simulate the sequencing, we generate a read by copying $\seq^\gmut$ letter by letter.
During copying, we may misread each letter, editing the \nth{i} letter of $\seq^\gmut$ with probability $\pphred{i}$, which depends on the letter's position $i$.

In \cref{app:model-equivalence}, we discuss the equivalence of the Core Genomes Model with the Graph Model.
We demonstrate that the distributions induced by both models are approximately the same, with different behavior only towards the end of the genome \todo{what happens here exactly?}.

\para{Including Mutations in the Graph Model}
The fourth column in \cref{fig:models} transforms the variant graph $H$ to a new variant graph $H'$ that directly generates mutated genome sequences.
This is desirable, as it means that our inference algorithm does not need to handle mutations explicitly, because we may handle mutations by incorporating them into the variant graph directly.

\cref{fig:enhance-variance-graph} shows this transformation, which replaces every edge $e$ in $H$ by $9$ edges ($4$ for insertions, $1$ for deletions, $3$ for edits, and $1$ for copy).
Insertions produce a new inserted letter without reading a letter from the core genome.
Therefore, $H'$ adds $4$ edges (each with probability $p \cdot \frac{\pins}{4}$) that allow emitting a letter without changing the current state.
Deletions allow reading a letter from the core genome without copying it.
Therefore, $H'$ adds an edge that does not emit a letter (it is labeled by $\epsilon$) but changes the state.
Edits allow reading a letter from the core genome but copying a different letter.
Therefore, $H'$ adds $3$ edges (each with probability $p \cdot \frac{\ped}{3}$) that emit a letter but change the state is if emitting the original letter.

\begin{figure}
	%\centering
	\small
	\begin{tikzpicture}[every text node part/.style={align=center}]
		\node[letter] at (-1.5,0) (v) {};
		\node[letter] at ( 1.5,0) (w) {};
		\draw[arrow,out=0,in=180+0] (v) edge node[above] {$A \; (p)$} (w);

		\node[inner sep=0pt,above right,anchor=center] at (0,-0.5) {\includegraphics[scale=1,angle=-90]{figures/arrow}};

		\node[letter] at (-1.5,-2) (v') {};
		\node[letter] at ( 1.5,-2) (w') {};
		% copy
		\draw[arrow,out=0,in=180+0] (v') edge node[above] {$A \; \left(p\cdot \pcopy\right)$} (w');
		% edit
		\draw[editarrow,out=30,in=180-30] (v') edge node[above] {$C/G/T \; \left( p\cdot \frac{\ped}{3} \right)$} (w');
		% delete
		\draw[editarrow,out=-30,in=180+30] (v') edge node[below] {$\epsilon \; \left(p \cdot \pdel \right)$} (w');
		% insert
		\draw[editarrow,in=180-40,out=180+40,looseness=12] (v') edge node[left] {$A/C/G/T$ \\ $\left( p \cdot \frac{\pins}{4} \right)$} (v');
	\end{tikzpicture}
	\caption{Enhancing $H$ to account for mutations. For illustration purposes, $H$ only consists of a single edge emitting letter $A$ with probability $p$. We list multiple letters to indicate multiple edges (\eg $A/C/G/T$).}
	\label{fig:enhance-variance-graph}
\end{figure}	